\chapter{Introduction}\label{sec:introduction}

\section{Spin glass dynamics and aging}\label{sec:intro-background}

The static theory of mean-field spin glasses describes the Gibbs measure of a random Hamiltonian in high dimension. For the spherical $p$-spin model studied in this paper, the limiting free energy is given by the Crisanti--Sommers formula, a variational problem over probability measures on $[0,1]$ whose minimizer is the Parisi measure \cite{cs1992}. This formula was proved by Talagrand for even models \cite{talagrand2006} and by Chen in general \cite{chen2013}. For the pure model, the Parisi measure is a point mass at zero at high temperature, and it has one step of replica symmetry breaking at every temperature below the static transition \cite[Section~5]{auffingerchen2015}; the two-point support property follows from \cite[Proposition~2.2]{talagrand2006}. These results concern the equilibrium measure. In this paper, we study the dynamics.

The natural dynamics is the Langevin diffusion on the sphere, in which the configuration moves in the direction of increasing Hamiltonian, perturbed by Brownian noise; the inverse temperature sets the strength of the drift relative to the noise. This diffusion is reversible with respect to the Gibbs measure, and for each fixed dimension $N$, it converges to the Gibbs measure as time tends to infinity. At sufficiently low temperature, however, the spectral gap of its generator is exponentially small in $N$, with probability tending to one over the disorder \cite{gj2019}, so the slowest relaxation modes have relaxation times that grow exponentially with $N$. On time scales that do not grow with $N$, the dynamics can stay far from equilibrium. In this paper, we take the limit $N\to\infty$ first, on every fixed time interval, and only then study the long-time behavior of the limiting dynamics. In this order of limits, the Gibbs measure and its Parisi measure need not describe the behavior at long times. The physical theory predicts instead that the system ages.

Aging is observed in disordered materials that do not reach equilibrium on any accessible time scale. After a rapid cooling to low temperature, called a quench, the response of such a material to a small perturbation depends on the time that has elapsed since the quench, and it becomes slower as the sample gets older. In spin glasses, Lundgren, Svedlindh, Nordblad, and Beckman found that magnetic relaxation in CuMn depends on the time spent at low temperature before the measurement \cite{lundgren1983}; the experimental picture is reviewed in \cite{vincent1997}. In a trap model, the system jumps between traps and spends a random time in each of them. Bouchaud introduced weak ergodicity breaking through a trap model whose mean trapping time diverges: the system keeps aging, although the phase space does not split into mutually inaccessible regions \cite[p.~1707]{bouchaud1992}. Here the correlation measures the overlap of the configurations at two times, and the response measures the linear effect of a small force applied at the earlier time; we define both precisely in \cref{sec:intro-model}. In the formulation used for mean-field aging, the correlation with every fixed earlier configuration eventually vanishes, while relaxation becomes slower as the system gets older \cite[arXiv version~1, Section~3]{ck1994}.

Mean-field spin glasses with $p$-spin interactions are a setting in which aging can be derived from a microscopic dynamics. Derrida introduced the random energy model, in which the energies of different configurations are independent, as a limit of a family of $p$-spin models \cite{derrida1980}, and the Ising $p$-spin models themselves were studied by Gross and M\'ezard \cite{grossmezard1984} and by Gardner \cite{gardner1985}. Sompolinsky and Zippelius developed the dynamical mean-field theory of spin glass relaxation \cite{sz1981,sz1982}. This theory replaces the interaction of one coordinate with all the others by a random force and a memory of its own past, both determined self-consistently, and it describes the dynamics through closed equations for the correlation and response functions. Kirkpatrick and Thirumalai connected the dynamics of $p$-spin models with the glass transition of supercooled liquids, which are liquids cooled below their freezing temperature without crystallizing \cite{kt1987a,kt1987b}. For $p\ge3$, the dynamical equations resemble those of the mode-coupling theory of supercooled liquids, a self-consistent theory in which correlation functions satisfy nonlinear equations with memory and, at low temperature, develop a plateau, that is, a positive limit at large time differences (see \cite{goetzesjogren1992} for a review). In these equations, the correlation of the equilibrium dynamics stops decaying to zero below a dynamical transition temperature. This temperature is higher than the static transition temperature, at which the Parisi measure changes. Kirkpatrick, Thirumalai, and Wolynes built on this correspondence a scaling theory of glasses in finite dimension, in which relaxation below the dynamical transition proceeds by thermally activated rearrangements, where the noise carries the system across barriers of the energy landscape \cite{ktw1989}. The spherical $p$-spin model was solved at equilibrium by Crisanti and Sommers \cite{cs1992}. Crisanti, Horner, and Sommers studied its dynamics in \cite{chs1993}, including the dynamical transition above the static one and the power-law relaxation near it. Cugliandolo and Kurchan constructed an asymptotic solution of the nonequilibrium equations at long times \cite{ck1993}. We describe their predictions in \cref{sec:intro-picture}, after we introduce the model and the equations that they analyzed.

\section{The model and the Cugliandolo--Kurchan equations}\label{sec:intro-model}

We fix an integer $p\ge3$ and $\beta>0$, and write $\xi(q)=q^p$. For $N\ge2$, let $S_N=\{x\in\R^N:\norm{x}^2=N\}$. The Hamiltonian $H_N$ is a centered Gaussian process on $\R^N$ with covariance
\[
 \E H_N(x)H_N(y)=N\xi\Bigl(\frac{x\cdot y}{N}\Bigr),
\]
realized as in \eqref{eq:hamiltonian}. This is the pure spherical $p$-spin model in the normalization of the rigorous static theory. Its Gibbs measure at inverse temperature $\beta$ has density proportional to $e^{\beta H_N}$ with respect to the uniform probability measure on $S_N$. Let $X_t$ be the Langevin dynamics on $S_N$ with generator
\[
 \Delta_S+\beta\nabla_SH_N\cdot\nabla_S,
\]
where $\Delta_S$ and $\nabla_S$ are the Laplace--Beltrami operator and gradient on $S_N$; equivalently, $X_t$ solves \eqref{eq:hard-sde}, whose noise has amplitude $\sqrt2$. The Gibbs measure is reversible for this diffusion. The initial point $X_0$ is uniform on $S_N$ and independent of the disorder and of the Brownian motion. Since the uniform measure is the Gibbs measure at $\beta=0$, this initial condition models a quench from infinite temperature. The drift points in the direction of increasing $H_N$, so the limiting energies and the threshold energy in this paper are positive; physics works usually minimize the Hamiltonian, which reverses these signs.

Since the uniform initial law is not invariant for the dynamics, its natural observables depend on two times. The correlation
\[
 C_N(t,s)=\frac{X_t\cdot X_s}N
\]
is the overlap of the configurations at times $s$ and $t$. The response describes how the configuration at a time $t$ reacts to a small force applied at an earlier time $s$. Precisely, for a deterministic force $h(u)$ that is bounded and continuous on each fixed time horizon and projected onto the tangent space of $S_N$, the response is the zero-initial strong solution $V_t$ of the forced linearized equation along $X$, with the drift and noise derivatives evaluated on $X$. It satisfies $V_t=\int_0^tU_N(t,u)h(u)\,\dd u$ for the random response matrix $U_N(t,u)$ (\cref{sec:langevin}). We write
\[
 \bar r_N(t,s)=\frac1N\E\tr U_N(t,s),\qquad 0\le s\le t,
\]
for its normalized trace, averaged over the disorder, the initial point, and the noise. Thus $\bar r_N(t,s)$ is the mean linear response of a coordinate at time $t$ to a unit impulse of force on the same coordinate at time $s$, that is, to a force concentrated near time $s$ with unit time integral. We measure the response to a force added to the drift without a factor $\beta$. The finite-dimensional response has diagonal value $\bar r_N(t,t)=\kappa_N=(N-1)/N$, because the force is projected onto the $(N-1)$-dimensional tangent space, and its limit below has diagonal value one (\cref{rem:fdt-convention}).

As $N\to\infty$, these functions have deterministic limits on fixed time horizons. For every fixed $0\le T<\infty$, the correlation $C_N(t,s)$ converges in probability to a function $C(t,s)$, uniformly on $[0,T]^2$, and the averaged response $\bar r_N(t,s)$ converges to a function $R(t,s)$, uniformly on $0\le s\le t\le T$. Together with a function $\lambda(t)$, which keeps $C(t,t)=1$, the limits form the unique solution of the system of integro-differential equations
\begin{align*}
 \partial_tR(t,s)&=-\lambda(t)R(t,s)
 +\beta^2\int_s^t\xi''(C(t,u))R(t,u)R(u,s)\,\dd u,\\
 \partial_tC(t,s)&=-\lambda(t)C(t,s)
 +\beta^2\int_0^t\xi''(C(t,u))R(t,u)C(u,s)\,\dd u\\
 &\hspace{1em}+\beta^2\int_0^s\xi'(C(t,u))R(s,u)\,\dd u,\\
 \lambda(t)&=1+\beta^2\int_0^t
 \bigl[\xi'(C(t,u))+C(t,u)\xi''(C(t,u))\bigr]R(t,u)\,\dd u,
\end{align*}
for $t\ge s\ge0$, with $C(t,t)=R(t,t)=1$, where $C$ is extended to $s>t$ by symmetry. We call these the Cugliandolo--Kurchan (CK) equations; they are stated with their conventions in \eqref{eq:ck-response}--\eqref{eq:ck-boundary}. Existence and uniqueness of the solution were proved by Dembo, Guionnet, and Mazza \cite[Proposition~1.1]{dgm2007}, and the solution satisfies $0\le C\le1$ and $R\ge0$ (see \cref{prop:ck-wellposed}). Sellke proved fixed-horizon convergence of $C_N$ and of the Brownian observable representing the integrated response \cite[Theorem~4.3]{sellke2024}. The force-response identification in \cref{lem:brown-force-response} equates the annealed mean of its rescaled version $\Gamma_N(t,s)=X_t\cdot B_s/(\sqrt2\,N)$ with $\int_0^s\bar r_N(t,u)\,\dd u$, where $B$ drives the diffusion. He built on the theorem of Ben Arous, Dembo, and Guionnet for the soft-spherical dynamics \cite[Theorem~1.2]{bag2006}, in which the constraint $\norm{X_t}^2=N$ is replaced by a confining potential. We state these results in our normalization in \cref{prop:hard-bridge,lem:brown-force-response}, and we deduce the uniform convergence of $\bar r_N$ in \cref{rem:no-pointwise-response}. Throughout the paper, the \emph{CK solution} is this solution $(C,R,\lambda)$, which describes the Langevin dynamics started from the uniform point.

The CK equations involve only $C$, $R$, and $\lambda$; the disorder no longer appears. The function $\lambda$ is a Lagrange multiplier for the spherical constraint, and the integrals are memory terms. In the physical derivation through dynamical mean-field theory, each coordinate of $X_t$ behaves in the limit like the solution $x_t$ of the linear equation (see \cite[arXiv version~2, Section~9.5]{cugliandolo2004})
\[
 \dd x_t=\Bigl(-\lambda(t)x_t+\beta^2\int_0^t\xi''(C(t,u))R(t,u)x_u\,\dd u+\eta_t\Bigr)\dd t+\sqrt2\,\dd b_t.
\]
Here $x_0$, $b$, and $\eta$ are independent; $x_0$ is a standard Gaussian variable, $b$ is a standard Brownian motion, and $\eta$ is a centered Gaussian process with covariance $\E\eta_t\eta_u=\beta^2\xi'(C(t,u))$. The interaction of one coordinate with all the others produces the memory kernel $\beta^2\xi''(C(t,u))R(t,u)$, which weights the effect of the past values of $x$ on its present drift, and the noise $\eta$, which is called colored because its values at different times are correlated. The CK equations express the self-consistency conditions that $C(t,s)=\E x_tx_s$, that $R(t,s)$ is the response of $x_t$ to a force applied at time $s$, and that $\E x_t^2=1$, which determines $\lambda$. We use this effective process only as a guide. Its variation-of-constants formula writes $C(t,s)$ as the sum of the contributions of the initial condition, of the Brownian noise, and of the colored noise, each propagated by the response. This is the Green covariance identity of \cref{prop:green-covariance}, which we prove directly from the CK equations.

We write
\[
 e_\beta(t)=\beta p\int_0^tC(t,u)^{p-1}R(t,u)\,\dd u,\qquad
 \chi(t,s)=\int_s^tR(t,u)\,\dd u.
\]
The first quantity is the limit of the mean energy $N^{-1}\E H_N(X_t)$, and $\lambda(t)=1+p\beta e_\beta(t)$ (\cref{prop:energy-bridge}). The second is the integrated response, or susceptibility. It is the response at time $t$ to a constant unit force switched on at time $s$. A field $h$ added to the Hamiltonian adds the force $\beta h$ to the drift, so the magnetization per unit field is $\beta\chi(t,s)$ in our conventions. It is the analog of the magnetization measured in aging experiments after a small field is switched on at time $s$, the zero-field-cooled protocol \cite{vincent1997}.

We also use the entropy production of the dynamics. Fix the disorder, and let $\rho_t$ be the law of $X_t$. Since the dynamics is reversible, the relative entropy of $\rho_t$ with respect to the Gibbs measure decreases in time. Its rate of decrease is the relative Fisher information, the integral against $\rho_t$ of the squared norm of the gradient of the logarithm of the density of $\rho_t$ with respect to the Gibbs measure. We call this rate the entropy production at time $t$. The initial relative entropy is at most $\beta^2N/2$ on average over the disorder, so the total entropy production per coordinate is at most $\beta^2/2$ on average (\cref{lem:entropy-budget}). We write $a_N(t)$ for the entropy production at time $t$, divided by $N$ and averaged over the disorder. Following the usual terminology for disorder averages, we also call $a_N(t)$ the normalized annealed Fisher information (\cref{sec:entropy-clusters}).

Other works use other normalizations. Cugliandolo and Kurchan \cite[Equations~(1)--(2)]{ck1993} minimize a Hamiltonian whose couplings are indexed by distinct unordered $p$-tuples and have variance $p!/(2N^{p-1})$. Their limiting covariance convention is $N\xi(x\cdot y/N)/2$, with an inverse temperature $\beta^{\mathrm{CK}}$ in front of the drift and noise with variance rate $2$. At the level of the limiting equations, their convention agrees with ours when $\beta^{\mathrm{CK}}=\sqrt2\,\beta$, with no change of time. In the physical normalization of time, where the drift is the gradient itself and the noise has variance rate $2\theta$ at temperature $\theta$, a rescaling of time gives $\beta^{\mathrm{CK}}=1/\theta$, and hence $\beta=1/(\sqrt2\,\theta)$. The mathematical works \cite{bag2006,dgm2007,sellke2024} use Brownian noise with variance rate one; their drift parameter is $\alpha=\beta/2$, and their time is $v=2t$. \Cref{tab:conventions} in \cref{sec:conventions} collects the dictionary. All formulas in this paper are written in our conventions, and the formulas that we quote from \cite{ck1993,cs1995,dgm2007,sellke2024} have been translated.

\section{The physical picture}\label{sec:intro-picture}

We first recall what happens in equilibrium. Suppose that the diffusion is started from the Gibbs measure instead of the uniform measure. Then the process is stationary, and its correlation and response depend only on the time difference $\tau=t-s$, which we also call the lag: they have the form $c(\tau)$ and $r(\tau)$. They satisfy the fluctuation--dissipation theorem (FDT),
\[
 r(\tau)=-c'(\tau),\qquad \tau>0,
\]
which expresses the response to a small force through the spontaneous fluctuations of the system. FDT is a consequence of reversibility (\cref{rem:fdt-convention}); see Kubo \cite{kubo1957} for the general linear-response theory. Agarwal \cite{agarwal1972} treats response around nonthermal steady states under detailed balance; nonequilibrium diffusion response identities are discussed in \cite{bmw2009}. In our conventions it carries no factor $\beta$. When $C(t,s)=c(t-s)$, it reads $R(t,s)=\partial_sC(t,s)$. Reversibility also gives a spectral representation. By the spectral theorem for the generator, the equilibrium correlation is a mixture of decaying exponentials, $c(\tau)=\int_{[0,\infty)}e^{-x\tau}\,\nu(\dd x)$ for a probability measure $\nu$. By Bernstein's theorem \cite{widder1941}, a continuous function $c$ on $[0,\infty)$ with $c(0)=1$ has this form if and only if it is \emph{completely monotone} on $(0,\infty)$, that is, infinitely differentiable there with
\[
 (-1)^nc^{(n)}(\tau)\ge0,\qquad n\ge0,\ \tau>0.
\]
In particular, a completely monotone function is nonnegative, nonincreasing, and convex. We use the terms FDT and complete monotonicity in this sense throughout.

At high temperature, the dynamics started from the uniform point approaches this equilibrium behavior. Dembo, Guionnet, and Mazza proved that for $\beta$ small enough, the CK solution relaxes exponentially fast and converges to a stationary solution that satisfies FDT \cite[Proposition~1.2 and Theorem~1.3]{dgm2007}.

Below the dynamical transition, the picture is different. Let $s$ be the waiting time, which is the age of the system when a measurement starts, and let $\tau=t-s$ be the lag. Cugliandolo and Kurchan constructed an asymptotic solution of the CK equations at long times \cite{ck1993}. It has the following features; the equation numbers refer to the arXiv version of \cite{ck1993}.
\begin{itemize}
\item \emph{A stationary regime.} When $s\to\infty$ at fixed lag $\tau$, we have $C(s+\tau,s)\to c(\tau)$ and $R(s+\tau,s)\to r(\tau)$, and the limits satisfy FDT. On bounded windows of late times, the correlation and the response behave as in equilibrium.
\item \emph{A positive plateau.} The stationary correlation does not decay to zero. It decreases to a positive limit $q=\lim_{\tau\to\infty}c(\tau)$, called the plateau; in the physics literature, $q$ is also called the dynamical Edwards--Anderson parameter.
\item \emph{Weak ergodicity breaking and an aging regime.} Below the plateau, the correlation decays on time scales that grow with the waiting time $s$, and $C(t,s)\to0$ as $t\to\infty$ for every fixed $s$. The aging regime consists of pairs of late times $s<t$ with $t-s\to\infty$ at which $C(t,s)$ stays below the plateau.
\item \emph{Marginality.} The plateau is fixed by a marginality condition \cite[Equation~(7)]{ck1993}, which says that a stability eigenvalue vanishes; we explain its meaning below. In our conventions, it reads $\beta^2p(p-1)q^{p-2}(1-q)^2=1$.
\item \emph{A constant fluctuation--dissipation ratio.} In the aging regime, FDT fails, and response and correlation are related instead by $R(t,s)\approx x\,\partial_sC(t,s)$, with a constant $x=(p-2)(1-q)/q$ given in the paragraph after \cite[Equation~(12)]{ck1993}; the corresponding response moments are \cite[Equation~(13)]{ck1993}. Equilibrium FDT corresponds to $x=1$, and below the dynamical transition $x<1$. Cugliandolo, Kurchan, and Peliti later interpreted $x$ as the ratio of the temperature of the noise to an effective temperature of the slow motion, which is therefore higher \cite{ckp1997}.
\item \emph{The energy.} The limit of the energy $e_\beta(t)$ as $t\to\infty$ is given by \cite[Equation~(9)]{ck1993}, combined with the relation between response and correlation in the aging regime; in our conventions, it is $\beta[1+(p-2)q^{p-1}-(p-1)q^p]$.
\end{itemize}
The first three features imply that the two iterated limits of the correlation differ:
\[
 \lim_{\tau\to\infty}\lim_{s\to\infty}C(s+\tau,s)=q>0,
 \qquad
 \lim_{s\to\infty}\lim_{\tau\to\infty}C(s+\tau,s)=0.
\]
Taking the limit of large waiting time first describes the relaxation of an arbitrarily old system over a bounded time window. Taking the limit of large lag first asks whether the configuration at a fixed time is eventually forgotten. In equilibrium, both limits equal $\lim_{\tau\to\infty}c(\tau)$. In this model, their difference is a mathematical signature of aging: at large waiting times, the correlation first approaches the plateau, as it would in equilibrium, and it then decays further on a time scale that diverges with the waiting time. Thus stationarity on bounded windows does not describe the relaxation over the whole history. The noncommutation of the limits detects this slower relaxation without determining its time scale.

The dynamical transition is the temperature at which $x$ reaches one \cite{ck1993}. In our conventions, the marginality condition and $x=1$ give the transition parameters
\[
 \qd=\frac{p-2}{p-1},\qquad \bd^2=\frac{(p-1)^{p-1}}{p(p-2)^{p-2}}.
\]
For the cubic model, $\qd=1/2$ and $\bd=\sqrt{4/3}$. For $\beta<\bd$, the physical picture predicts a zero plateau \cite{ck1993,chs1993}. Equivalently, $\bd$ is the smallest $\beta$ for which the equation $\beta^2pq^{p-2}(1-q)=1$ has a solution $q\in(0,1)$, and this solution is $\qd$ (\cref{sec:conventions}). In the physical picture, this is the equation for the plateau of the equilibrium dynamics in the replica symmetric phase, so the equilibrium dynamics stops decorrelating at the inverse temperature $\bd$ \cite{chs1993}.

The picture can be compared with the static theory. Let $\beta_{\mathrm s}$ be the static transition, so that the Parisi measure is the point mass at zero for $\beta\le\beta_{\mathrm s}$ and has two atoms for $\beta>\beta_{\mathrm s}$ \cite{auffingerchen2015}. Then $\bd<\beta_{\mathrm s}$; for $p=3$, we have $\bd\approx1.155$ and $\beta_{\mathrm s}\approx1.207$ (see \cite[Theorem~1]{subag2023} for the equations that determine $\beta_{\mathrm s}$ in our normalization). Since $\xi''(q)=p(p-1)q^{p-2}$, the marginality condition reads $\beta^2\xi''(q)(1-q)^2=1$. Consider a trial measure with one step of replica symmetry breaking, $\theta\delta_0+(1-\theta)\delta_q$ with $\theta\in(0,1)$, in the Crisanti--Sommers functional. In the replica calculation, its stability against perturbations of the overlaps within a block is measured by the replicon eigenvalue, which is a positive multiple of $1-\beta^2\xi''(q)(1-q)^2$. Thus the marginality condition says that the replicon eigenvalue vanishes, so that the solution is marginally stable. If, in addition, the derivative of the functional in $q$ vanishes, then $\theta=x$. For $\beta>\beta_{\mathrm s}$, the Parisi measure is obtained instead by optimizing the functional over both $\theta$ and $q$, and it is strictly stable. Barrat, Franz, and Parisi identified the asymptotic dynamical condition with this static marginal stability \cite[Section~IV, Equation~(25)]{bfp1997}; see also \cite[arXiv version~2, Section~15.1]{cugliandolo2004}.

In this sense, the plateau and the energy of the dynamics are those of marginally stable states rather than those of the equilibrium states that the Parisi measure describes, and the fluctuation--dissipation ratio $x$ is the weight $\theta$ of the corresponding marginal solution. Cugliandolo and Kurchan noted that, in this model, the overlap distribution defined by the dynamics differs from the static one \cite[the paragraph after Equation~(15)]{ck1993}. Franz, M\'ezard, Parisi, and Peliti related the dynamical response to the static Parisi measure under additional hypotheses, and they pointed out that the pure spherical $p$-spin model does not satisfy this relation \cite{fmpp1998}. For $\bd<\beta\le\beta_{\mathrm s}$, the Parisi measure carries no positive overlap, while the predicted dynamical plateau is positive.

The limiting energy also has a landscape interpretation, developed alongside the calculation of Cugliandolo and Kurchan. The Thouless--Anderson--Palmer (TAP) free energy is a function of the local magnetizations whose stable local minima represent candidate metastable states. Crisanti and Sommers counted the solutions of the TAP equations of the spherical model and computed the threshold energy at which these solutions become marginally stable \cite[Equations~(8)--(9)]{cs1995}. At this threshold, the lower edge of the limiting spectrum of the Hessian of the TAP free energy reaches zero \cite[arXiv version~1, p.~4]{ck1993}; see also \cite[arXiv version~2, Section~15.4]{cugliandolo2004}. The energy of these threshold states depends on the temperature, and it agrees with the limiting energy of Cugliandolo and Kurchan \cite[Equation~(14)]{cs1995}. As $\beta\to\infty$, it tends to the threshold energy $\Eth=2\sqrt{(p-1)/p}$ of the landscape, which is strictly smaller than the ground-state energy $\lim_{N\to\infty}N^{-1}\E\max_{S_N}H_N$; for $p=3$, we have $\Eth\approx1.633$, while the ground-state energy is approximately $1.657$ \cite{chensen2017}. Kurchan and Laloux related slow dynamics to the geometry of the landscape near the threshold \cite[Section~V]{kurchanlaloux1996}. Related analyses of metastable states, barriers, and stationary points of the TAP free energy are in \cite{kpv1993,cgp1998}.

Cugliandolo and Kurchan developed the structure of the aging regime further in \cite{ck1994}, where they studied the Sherrington--Kirkpatrick model. They distinguished weak ergodicity breaking from weak long-term memory, which means that the integrated response to a force applied during any fixed time interval $[0,S]$ vanishes as $t\to\infty$ \cite[arXiv version~1, Section~3]{ck1994}. They also observed that the asymptotic equations of the aging regime do not change when time is replaced by an increasing function of time and the response is transformed accordingly, so that these equations do not determine the time scale of aging \cite[arXiv version~1, Section~4]{ck1994}. Further, they studied the relations between correlations at three widely separated times \cite[arXiv version~1, Section~5]{ck1994}. Two important examples of such relations are the dynamical ultrametric relation $C(t,u)=\min\{C(t,s),C(s,u)\}$ for $u<s<t$, analogous to the ultrametric organization of equilibrium overlaps, and a product relation; for spherical $p$-spin models, they discussed a product relation \cite[arXiv version~1, Section~5.1, Equation~(5.1.8)]{ck1994}. Reviews of the two regimes and their response are given in \cite{bckm1998,cugliandolo2004}. The idealized mode-coupling scenario of the glass transition was developed by Bengtzelius, G\"otze, and Sj\"olander and by Leutheusser \cite{bgs1984,leutheusser1984}. Mode-coupling theory also describes how a critical or marginal stationary correlation $c(\tau)$ approaches its plateau: $c(\tau)-q$ decays like $\tau^{-a}$, where $a\in(0,1/2)$ solves $\Gamma(1-a)^2/\Gamma(1-2a)=\lambda_{\mathrm{MCT}}$, $\Gamma$ is Euler's gamma function, and the constant $\lambda_{\mathrm{MCT}}$, called the exponent parameter, is not to be confused with the multiplier $\lambda(t)$. The correspondence between disordered models and mode-coupling theory was worked out by Bouchaud, Cugliandolo, Kurchan, and M\'ezard in \cite{bckm1996}. They gave this exponent relation at equilibrium \cite[arXiv version, Equation~(IV.3)]{bckm1996} and for the stationary profile of the aging solution \cite[arXiv version, Equation~(IV.6)]{bckm1996}, where they attribute the calculation to Cugliandolo and Le Doussal \cite{cugliandololedoussal1996}. For the pure spherical model, the exponent parameter of the aging solution is $x/2=(p-2)(1-q)/(2q)$. This relation is written explicitly by Kim and Latz \cite[the paragraph after Equation~(12)]{kimlatz2001}; see also \cite[arXiv version~2, Equation~(12.36)]{cugliandolo2004}.

The order of limits is essential in this picture. As explained in \cref{sec:intro-background}, the dynamics converges to its Gibbs measure for each finite $N$. In the CK description, the dimension tends to infinity first, on every fixed time interval, and only then do the observation and waiting times tend to infinity. Aging in this sense is a property of a deterministic pair of functions $(C,R)$, and it is not a statement about time scales that grow with $N$. On such time scales, the dynamics can be activated, and the relevant mechanisms are different; we discuss the rigorous results on activated aging in \cref{sec:intro-related}. Neither a theorem about finite-dimensional mixing nor a formal solution of the asymptotic equations determines the behavior of the CK solution in the order of limits used here.

\section{Main results}\label{sec:main-results}

The goal of this paper is to prove the main features of the picture of \cref{sec:intro-picture} for pure $p$-spin models with $p\ge3$ started from a uniform point. The physical calculation imposes a long-time form on the solution and solves reduced asymptotic equations. We prove its conclusions for the CK solution itself, which is the limit of the Langevin dynamics, without assuming that form. At every positive temperature, we prove that the CK solution becomes stationary on bounded windows of lags and satisfies FDT there, that its energy converges, and that its correlation and response with every fixed earlier time tend to zero. For $\beta\ge\bd$, that is, at and below the dynamical transition temperature, we prove that the plateau is marginal and that the limiting energy is the one predicted by Cugliandolo and Kurchan. The two iterated limits of the correlation then differ, which is the sense in which the solution ages. We also prove that the response at a late time to forces applied at all earlier times concentrates at short lags in the mean-square sense, that the violation of FDT vanishes in a mean-square sense, and that, for $\beta>\bd$, its positive part keeps a positive integrated mass. For $\beta\ge\bd$, we prove that the stationary correlation approaches its plateau as a power law with a positive and finite amplitude. For the cubic model, we prove a lower bound on the susceptibility and an inequality between correlations at three times for $\beta\ge\bd$, and, for $\beta>\bd$, a lower bound that relates the time needed to relax below the plateau to the entropy production. We do not identify the time scale of aging, and we do not prove the full fluctuation--dissipation relation of the aging regime.

Let $(C,R,\lambda)$ be the CK solution of \cref{sec:intro-model}. All long-time limits in the following statements concern this deterministic solution. They describe the Langevin dynamics only in the order of limits described in \cref{sec:intro-picture}: first $N\to\infty$ on a fixed time horizon, through \cref{prop:hard-bridge}, and then the long-time limits.

\begin{theorem}[Stationary limit, marginal plateau, and aging]\label{thm:main-aging}
Fix $p\ge3$ and $\beta>0$. There exist functions $c_\beta,r_\beta\colon[0,\infty)\to\R$ and numbers $e_\beta(\infty)$ and $\ell_\beta$ such that $e_\beta(T)\to e_\beta(\infty)$ as $T\to\infty$ and
\begin{equation}
 C(T+t,T+s)\longrightarrow c_\beta(t-s),\qquad
 R(T+t,T+s)\longrightarrow r_\beta(t-s),\qquad
 \lambda(T)\longrightarrow\ell_\beta,
 \label{eq:main-stationary}
\end{equation}
where the first two limits are locally uniform in $t\ge s\ge0$. The function $c_\beta$ is completely monotone on $(0,\infty)$, $r_\beta=-c_\beta'$, and $c_\beta(0)=r_\beta(0)=1$. For every finite $S\ge0$,
\begin{equation}
 \lim_{t\to\infty}\sup_{0\le s\le S}\bigl\{C(t,s)+R(t,s)\bigr\}=0.
 \label{eq:main-forgetting}
\end{equation}
If $\beta\ge\bd$, then the plateau $q_\beta=\lim_{\tau\to\infty}c_\beta(\tau)$ is the larger of the two solutions $q\in(0,1)$ of
\begin{equation}
 \beta^2p(p-1)q^{p-2}(1-q)^2=1,
 \label{eq:main-marginality}
\end{equation}
and $q_\beta\ge\qd$. In this range,
\begin{align}
 e_\beta(\infty)
 &=\beta\bigl[1+(p-2)q_\beta^{p-1}-(p-1)q_\beta^p\bigr],
 \label{eq:main-energy}\\
 \ell_\beta
 &=\frac1{1-q_\beta}+\beta^2p\bigl(1-q_\beta^{p-1}\bigr),
 \label{eq:main-multiplier}
\end{align}
and the correlation limits do not commute:
\begin{equation}
 \lim_{\tau\to\infty}\lim_{s\to\infty}C(s+\tau,s)=q_\beta>0,
 \qquad
 \lim_{s\to\infty}\lim_{\tau\to\infty}C(s+\tau,s)=0.
 \label{eq:main-aging}
\end{equation}
\end{theorem}

\Cref{thm:main-aging} has two parts. The first holds at every temperature. It says that the CK solution becomes stationary at bounded lags: near every late time $T$, the correlation and the response are close to functions of the lag alone, these functions satisfy FDT, and the correlation has the spectral representation of an equilibrium correlation. Thus, on bounded windows of late times, the correlation and the response have the stationary FDT structure of an equilibrium dynamics, even at low temperature, where the dynamics keeps aging. The limits in \eqref{eq:main-stationary} are taken along the whole family $T\to\infty$, so the stationary profile does not depend on a subsequence. By \eqref{eq:main-forgetting}, the solution also forgets every fixed earlier time: its correlation with the configuration at a fixed time $s$, and its response to a force applied at time $s$, tend to zero.

The second part concerns $\beta\ge\bd$, and it identifies the plateau and the energy with those of the asymptotic solution of Cugliandolo and Kurchan. The relation \eqref{eq:main-marginality} is their marginality condition \cite[Equation~(7)]{ck1993}, and \eqref{eq:main-energy} is their energy \cite[Equation~(9)]{ck1993}, combined with the relation between response and correlation in the aging regime stated after \cite[Equation~(12)]{ck1993}; both are written here in the conventions of \cref{sec:conventions}. Since the plateau is positive while the correlation with every fixed earlier time vanishes, the two iterated limits in \eqref{eq:main-aging} differ. We call this noncommutation \emph{correlation aging}. It does not identify the time scale on which $C(s+\tau,s)$ decreases below $q_\beta$. In physical terms, \cref{thm:main-aging} shows that, for $\beta\ge\bd$, the dynamics started from a uniform point does not reach equilibrium: its plateau is the marginal one, and its limiting energy is the energy predicted for the threshold states. For example, if $\bd<\beta<\beta_{\mathrm s}$, the Parisi measure is the point mass at zero and the equilibrium energy per coordinate is $\beta$, while by \eqref{eq:main-energy}, we have $e_\beta(\infty)=\beta-\beta q_\beta^{p-1}[(p-1)q_\beta-(p-2)]<\beta$, since $q_\beta>\qd$.

The theorem includes the transition point $\beta=\bd$, where $q_\beta=\qd$ and $e_\beta(\infty)=\bd$. The value $e_{\bd}(\infty)=\bd$ coincides with the value $e_\beta(\infty)=\beta$ obtained whenever the plateau is zero (\cref{sel:intermediate-range}). Letting $\beta\to\infty$ in \eqref{eq:main-energy}, we find that $e_\beta(\infty)\to\Eth=2\sqrt{(p-1)/p}$ (\cref{rem:low-temperature}). This limit is the threshold energy of the pure $p$-spin landscape \cite[Equation~(2.13)]{abc2013}, and it agrees with the low-temperature results of Sellke for Langevin dynamics on bounded time horizons \cite[Theorems~1.1 and~1.7]{sellke2024}. The limit $\beta\to\infty$ is taken after all other limits, and it says nothing about gradient flow at zero temperature.

We prove \cref{thm:main-aging} at the end of \cref{sec:aging}, by combining \cref{sel:all-temperature,sel:selected-sector,sel:fixed-wait-forgetting,sel:aging} with \cref{prop:stationary-clusters,sel:cluster-complete-monotonicity}. The main difficulty is that the correlation equation involves the whole history: on late windows of bounded length, its memory integrals involve the overlap of late states with the remote past, which a limit on such windows does not control. The selection of the plateau and the forgetting of earlier times rest on two facts, which we state as \cref{prop:threshold-inputs} in \cref{sec:threshold-inputs}. To state them, we write $C(t,s)=\ip{v_t}{v_s}$ for unit vectors $v_t$ in a deterministic Hilbert space whose components are indexed by trees, and we project $v_t$ onto the components whose leaves record the initial condition or noise up to an earlier time, while their internal interaction times may be later (\cref{sec:geometry-core}). The representation comes with an energy functional, the tree energy, which reproduces $e_\beta(t)$ on $v_t$. The first fact says that, for every stationary limit with positive plateau, the tree energy of these projections, after normalization, does not exceed the threshold energy $\Eth$ in the limit. The second fact says that if the whole translated family converges to one stationary profile whose plateau is positive and marginal, then the solution forgets every fixed earlier time. We prove that proposition for $p\in\{3,4\}$ in \cref{sec:tree-ceilings}, using a deterministic inequality in this Hilbert space, and for every $p\ge3$ in \cref{app:general-selection}, using the algorithmic threshold of Huang and Sellke \cite[Theorem~1 and Corollary~1.8]{hs2023}.

To describe what happens for $\beta<\bd$, we use the following elementary facts about the marginality equation. Let
\[
 q_{\mathrm m}=\frac{p-2}{p},\qquad
 \beta_{\mathrm m}^2=\frac{p^{p-1}}{4(p-1)(p-2)^{p-2}}.
\]
The function $q\mapsto p(p-1)q^{p-2}(1-q)^2$ attains its maximum $\beta_{\mathrm m}^{-2}$ on $[0,1]$ at $q_{\mathrm m}$, so the marginality equation $\beta^2p(p-1)q^{p-2}(1-q)^2=1$ has a solution in $(0,1)$ if and only if $\beta\ge\beta_{\mathrm m}$; the value $\beta_{\mathrm m}$ is not a physical transition. For $\beta\ge\beta_{\mathrm m}$, we write $q_+(\beta)$ for its solution in $[q_{\mathrm m},1)$. We have $\beta_{\mathrm m}<\bd$, $q_{\mathrm m}<\qd$, and $q_+(\bd)=\qd$, and for $\beta\ge\bd$ we write $q_\beta=q_+(\beta)$; see \cref{sec:marginality,sec:selected-sector}.

For $\beta<\bd$, \cref{thm:main-aging} does not determine the plateau on the whole interval. For $\beta<\beta_{\mathrm m}$, the plateau is zero. For $\beta_{\mathrm m}\le\beta<\bd$, the proof shows that it is either zero or $q_+(\beta)$, and it proves convergence in both cases (\cref{sel:all-temperature,sel:intermediate-range}). The physical picture predicts the zero plateau there \cite{ck1993,chs1993}. We return to this question in \cref{sec:discussion-below}.

In the remaining statements, when $p$ and $\beta$ are fixed, we abbreviate $c=c_\beta$, $r=r_\beta$, $q=\lim_{\tau\to\infty}c_\beta(\tau)$, and $\ell=\ell_\beta$. All these statements use the stationary profile selected in \cref{thm:main-aging}. The only exception in notation is \cref{thm:main-fisher-cost}, where $c$ denotes a level in $(0,q)$, as in \cref{sec:three-time,resp:fisher-chapter}.

The stationary limit in \cref{thm:main-aging} describes the response $R(t,s)$ at bounded lags $t-s$, and the forgetting statement describes it at bounded earlier times $s$. The next theorem concerns the response at a late time $t$ to forces applied at all earlier times $s\in[0,t]$ simultaneously. We write it as a function of the lag $v=t-s\in[0,t]$, so that $R(t,t-v)$ is the response at time $t$ to an impulse at time $t-v$. We also write
\[
 E(t,s)=\partial_sC(t,s)-R(t,s),\qquad t>s>0,
\]
for the violation of FDT in the earlier time $s$, and $z_+=\max\{z,0\}$ for $z\in\R$. In equilibrium, $E=0$. For $\beta>\bd$, the physical picture predicts $E(t,s)\approx(1-x)\,\partial_sC(t,s)\ge0$ in the aging regime, and this quantity is expected to be small there, because the correlation changes on a diverging time scale.

\begin{theorem}[Response at late times and the violation of FDT]\label{thm:main-response}
Fix $p\ge3$ and $\beta>0$. Then, as $t\to\infty$,
\begin{equation}
\begin{aligned}
 \mathbf1_{[0,t]}(v)R(t,t-v)&\longrightarrow r(v),\\
 \mathbf1_{[0,t]}(v)\,\partial_sC(t,s)\big|_{s=t-v}&\longrightarrow r(v)
\end{aligned}
 \qquad\text{strongly in }L^2(0,\infty),
 \label{eq:main-strong-rows}
\end{equation}
where the functions on the left are set equal to zero for $v>t$. Moreover,
\begin{equation}
 \lim_{t\to\infty}\sup_{0\le s\le t}|R(t,s)-r(t-s)|=0,
 \qquad \chi(t,0)=o(\sqrt t).
 \label{eq:main-uniform-response}
\end{equation}
Let $E(t,s)=\partial_sC(t,s)-R(t,s)$. There exists a nonnegative integrable function $a_F$ on $[0,\infty)$, obtained as an $L^1_{\mathrm{loc}}$ limit of the normalized annealed Fisher informations along one subsequence of dimensions, such that $\int_0^\infty a_F(s)\,\dd s\le\beta^2/2$, $|E(t,s)|^2\le a_F(s)$ for $t>s>0$, and $\sup_{s\ge T}a_F(s)\to0$ as $T\to\infty$. Further,
\begin{equation}
 \lim_{t\to\infty}\int_0^t|E(t,s)|^2\,\dd s=0.
 \label{eq:main-defect-l2}
\end{equation}
If $\beta>\bd$, let
\[
 q=q_\beta,\qquad x=\frac{(p-2)(1-q)}q\in(0,1),\qquad
 m_\infty=q\bigl(1-x^{1/p}\bigr)>0.
\]
Then
\begin{equation}
 \lim_{A\to\infty}\lim_{t\to\infty}
 \frac{p\int_0^{t-A}C(t,u)^{p-1}R(t,u)\,\dd u}
 {C(t,t-A)^p-C(t,0)^p}=x,
 \label{eq:main-weighted-ratio}
\end{equation}
while, for every fixed $S\ge0$,
\begin{equation}
 \lim_{L\to\infty}\liminf_{t\to\infty}
 \int_S^{t-L}E(t,u)_+\,\dd u\ge m_\infty,
 \qquad
 \liminf_{t\to\infty}t\int_0^tE(t,u)_+^2\,\dd u\ge m_\infty^2.
 \label{eq:main-defect-budget}
\end{equation}
\end{theorem}

The first half of \cref{thm:main-response} says that the response and the derivative of the correlation in the earlier time concentrate at short lags. The stationary limit in \cref{thm:main-aging} controls the response only on bounded windows of lags, while \eqref{eq:main-strong-rows} concerns all lags $v\in[0,t]$, an interval whose length tends to infinity. In particular, for large $A$, the squared response integrated over lags larger than $A$ is small at late times, and by \eqref{eq:main-uniform-response}, the response converges uniformly over all earlier times. This does not imply that the response integrated over these lags vanishes: the total susceptibility $\chi(t,0)$ grows more slowly than $\sqrt t$, but we do not prove that it stays bounded. The convergence of the second function in \eqref{eq:main-strong-rows} is a strong form of FDT in the earlier time: the violation $E(t,\cdot)$ tends to zero in $L^2(0,t)$.

The bound $|E|^2\le a_F$ explains why FDT holds in this sense. The function $a_F$ is a limit, along one subsequence of dimensions, of the normalized entropy production $a_N$ of \cref{sec:intro-model}, and its integral is at most $\beta^2/2$. The bound $|E(t,s)|^2\le a_F(s)$ says that the violation of FDT at an earlier time $s$ is controlled by the entropy produced at time $s$. Since $\sup_{s\ge T}a_F(s)\to0$ as $T\to\infty$, the violation is uniformly small when the earlier time $s$ is large. This is the mechanism of Cugliandolo, Dean, and Kurchan \cite{cdk1997}, which we make uniform in the dimension. We call $a_F$ the Fisher majorant. It need not be unique, but we choose one dimension subsequence and one majorant and use them throughout, including in \cref{thm:main-fisher-cost}.

FDT also holds in the later time, uniformly over all earlier times, for every $p\ge3$ and $\beta>0$:
\[
 \lim_{t\to\infty}\sup_{0\le s\le t}|\partial_tC(t,s)+r(t-s)|
 =\lim_{t\to\infty}\sup_{0\le s\le t}|\partial_tC(t,s)+R(t,s)|=0
\]
(\cref{resp:forward-derivative}). On the diagonal, $\partial_tC(t,t)=-1$ denotes the derivative from the causal side with the earlier time held fixed; $c'(0)=-1$ is a right derivative.

The second half of \cref{thm:main-response} shows that the violation of FDT nevertheless carries a positive integrated mass in the aging regime. For every fixed $S\ge0$, the integral of the positive part of $E(t,\cdot)$ over $[S,t-L]$ has lower limit at least $m_\infty$ when first $t\to\infty$ and then $L\to\infty$, as in \eqref{eq:main-defect-budget}, while the $L^2$ norm of $E(t,\cdot)$ tends to zero. These statements hold together because the violation is integrated over the growing interval $[S,t-L]$, whose length is of order $t$; the theorem does not say how the violation is distributed over this interval. By the second inequality in \eqref{eq:main-defect-budget}, the $L^2$ norm of $E_+$ cannot decay faster than $t^{-1/2}$. Cugliandolo, Dean, and Kurchan explained in \cite{cdk1997} why entropy production allows small pointwise violations of FDT to coexist with a nonzero integrated violation over an expanding time interval. In \cref{thm:main-response}, we prove such a coexistence for the CK solution. The ratio in \eqref{eq:main-weighted-ratio} identifies one moment of the slow response with the fluctuation--dissipation ratio $x$ of \cite{ck1993}. Indeed, if the response satisfied $R(t,u)=x\,\partial_uC(t,u)$ for $u\le t-A$, as the constant ratio of \cite{ck1993} predicts for large $A$, then the ratio would equal $x$, since $p\int_0^{t-A}C(t,u)^{p-1}\partial_uC(t,u)\,\dd u=C(t,t-A)^p-C(t,0)^p$. We prove this weighted identity, but not the pointwise relation.

The Fisher majorant is constructed in \cref{thm:fisher-clock}. The convergence in \eqref{eq:main-strong-rows} and the bound $\chi(t,0)=o(\sqrt t)$ are proved in \cref{resp:strong-row,resp:uniform-row,resp:backward-fdt}, the identity \eqref{eq:main-weighted-ratio} in \cref{resp:weighted-ratio}, and the bounds \eqref{eq:main-defect-budget} in \cref{resp:persistent-defect}.

The convergence over all earlier times also shows where the correlation at the plateau comes from. The Green covariance identity of \cref{prop:green-covariance} writes $C(t,s)$ as the sum of three terms: the contribution $R(t,0)R(s,0)$ of the initial condition, a thermal term, which is the contribution of the Brownian noise, and a memory term, which is the contribution of the colored noise of the effective process. For $t\ge s$, the thermal term and its stationary counterpart are
\[
 B(t,s)=2\int_0^sR(t,u)R(s,u)\,\dd u,
 \qquad b_{\rm th}(\tau)=2\int_0^\infty r(\tau+v)r(v)\,\dd v.
\]
For every $p\ge3$ and $\beta>0$, we show in \cref{resp:thermal} that
\begin{equation}
 \lim_{S\to\infty}\sup_{t\ge s\ge S}|B(t,s)-b_{\rm th}(t-s)|=0.
 \label{eq:main-thermal}
\end{equation}
The function $b_{\rm th}$ is completely monotone and tends to zero. Hence the thermal term vanishes on every late pair of times whose lag diverges. The initial term also vanishes, by \eqref{eq:main-forgetting}. Thus any correlation retained on such a pair is carried by the memory term, which comes from the interaction of each coordinate with the others.

For the cubic model, we prove two further predictions of the physical picture in one-sided form. Cugliandolo and Kurchan described the aging regime through the relations between correlations at three times \cite[arXiv version~1, Section~5]{ck1994}. For spherical models, they discussed a relation in which correlations below the plateau compose multiplicatively: for three times that are far apart, $C(t,u)/q$ is the product of $C(t,s)/q$ and $C(s,u)/q$ \cite[arXiv version~1, Section~5.1, Equation~(5.1.8)]{ck1994}.

\begin{theorem}[Three-time inequality and susceptibility in the cubic model]\label{thm:main-geometry}
Let $p=3$ and $\beta\ge\bd$. For every sequence of triples $u_n<s_n<t_n$ with $u_n\to\infty$, $s_n-u_n\to\infty$, and $t_n-s_n\to\infty$,
\begin{equation}
 \liminf_{n\to\infty}
 \bigl\{qC(t_n,u_n)-C(t_n,s_n)C(s_n,u_n)\bigr\}\ge0.
 \label{eq:main-product}
\end{equation}
The same holds for $qC(t_n,s_n)-C(t_n,u_n)C(s_n,u_n)$ and for $qC(s_n,u_n)-C(t_n,s_n)C(t_n,u_n)$. Further, under the same hypotheses,
\begin{equation}
 \liminf_{t\to\infty}\chi(t,0)\ge2(1-q).
 \label{eq:main-susceptibility}
\end{equation}
\end{theorem}

For the cubic model, $\bd=\sqrt{4/3}$ and $q\ge1/2$, with $q=1/2$ at $\beta=\bd$; all three product inequalities and the susceptibility bound include this point. The inequality \eqref{eq:main-product} is a one-sided version of the product relation of Cugliandolo and Kurchan for the CK solution, with no assumption on the time scale of aging. In words, after dividing correlations by $q$, the correlation across two successive long time intervals is asymptotically at least the product of the correlations across the two intervals. For a level $c\in(0,q)$, consider the last time after $s$ at which the correlation with the configuration at time $s$ is at least $c$. The inequality \eqref{eq:main-product} implies that these last times compose: for fixed levels $c_1,c_2\in(0,q)$ and a fixed level below $c_1c_2/q$, and for all sufficiently large $s$, the last time for the lower level is at least the last time for $c_2$, started from the last time for $c_1$ (\cref{cor:geometry-last-crossing}). These statements allow the correlation to return above a level before its last crossing, and \eqref{eq:main-product} does not assert a product equality.

The bound \eqref{eq:main-susceptibility} concerns the total response to a constant force switched on at time zero. If FDT held at all pairs of times, we would have $\chi(t,0)=1-C(t,0)\to1$. Of this amount, the stationary regime contributes $\int_0^\infty r(v)\,\dd v=1-q$, so \eqref{eq:main-susceptibility} forces at least another $1-q$ of integrated response outside every fixed window of lags. In the full physical prediction, including the absence of response mass at correlation zero, the aging regime contributes $xq$ instead of $q$, which gives $\chi(t,0)\to(1-q)+xq=(p-1)(1-q)$; for $p=3$, this is $2(1-q)$, which is smaller than one when $q>1/2$. We prove only the lower bound. Both results come from the geometry of the tree vectors of \cref{sec:geometry-core} whose energies approach the threshold: the product inequality from \cref{thm:geometry-threshold-product}, and the susceptibility bound from a limiting operator associated with such vectors (\cref{thm:geometry-universal-root}). They are proved in \cref{thm:geometry-actual-product,thm:geometry-susceptibility-lower}.

The next result relates relaxation below the plateau to entropy production. For a level $c\in(0,q)$ and a waiting time $s$, we consider the time $\tau_c(s)$ that elapses after $s$ until the correlation with the configuration at time $s$ first reaches the level $c$.

\begin{theorem}[Entropy production and relaxation below the plateau in the cubic model]\label{thm:main-fisher-cost}
Let $p=3$ and $\beta>\bd$, and let $a_F$ be the Fisher majorant of \cref{thm:main-response}. If $s_n\to\infty$, $\tau_n\to\infty$, and
\[
 \tau_n\int_{s_n}^{s_n+\tau_n}a_F(u)\,\dd u\longrightarrow0,
\]
then $C(s_n+\tau_n,s_n)\to q$ and $\chi(s_n+\tau_n,s_n)\to1-q$. For $0<c<q$, the first passage time
\[
 \tau_c(s)=\inf\{\tau\ge0:C(s+\tau,s)\le c\}
\]
is finite, tends to infinity as $s\to\infty$, and satisfies
\begin{equation}
 \liminf_{s\to\infty}\tau_c(s)
 \int_s^{s+\tau_c(s)}a_F(u)\,\dd u
 \ge\min\biggl\{1,
 \biggl[\frac{3\beta^2(2q-1)(q-c)^2}{\ell+15\beta^2}\biggr]^2\biggr\}>0.
 \label{eq:main-passage-cost}
\end{equation}
The same lower bound holds after replacing the integral by $\int_s^\infty a_F(u)\,\dd u$.
\end{theorem}

The first statement concerns windows $[s_n,s_n+\tau_n]$ whose starting times and lengths both tend to infinity. If the length of the window times the integral of $a_F$ over it tends to zero, then at the end of the window, the correlation with its start tends to the plateau, and the susceptibility tends to its stationary value $1-q=\int_0^\infty r(v)\,\dd v$, as if the solution were stationary. Relaxation below the plateau therefore requires entropy production: to reach a fixed level $c<q$, the elapsed time multiplied by the integral of $a_F$ over the elapsed time must stay bounded below, and \eqref{eq:main-passage-cost} gives an explicit lower bound. As in \cref{resp:fisher-chapter}, we call the quantity bounded in \eqref{eq:main-passage-cost} the Fisher cost of the passage. Denoting the right side of \eqref{eq:main-passage-cost} by $\kappa_c$, the bound with $\int_s^\infty a_F$ shows that
\[
 \tau_c(s)\ge\frac{\kappa_c-o(1)}{\int_s^\infty a_F(u)\,\dd u}\qquad(s\to\infty).
\]
Since $a_F$ is integrable, the right side tends to infinity: the time needed to relax below the plateau grows at least as fast as the reciprocal of the remaining integral $\int_s^\infty a_F$ of the chosen Fisher majorant. This is a quantitative form of aging, but it does not give a growth law for $\tau_c(s)$, because we do not know how fast $\int_s^\infty a_F$ decays. The strict condition $\beta>\bd$ is essential to this positive bound: its factor $2q-1$ vanishes at the transition. The proof, in \cref{resp:fisher-cost}, combines the backward response equation with the product inequality of \cref{thm:main-geometry}.

Our last main result says that, for $\beta\ge\bd$, the stationary correlation approaches its plateau algebraically: $c(t)-q$ is asymptotic to a positive constant times $t^{-a}$, with the exponent predicted by mode-coupling theory.

\begin{theorem}[Power-law approach to the plateau]\label{thm:main-stationary-tail}
Fix $p\ge3$ and $\beta\ge\bd$. There exists a unique $a\in(0,1/2)$ such that
\begin{equation}
 \frac{\Gamma(1-a)^2}{\Gamma(1-2a)}=\frac{(p-2)(1-q)}{2q},
 \label{eq:main-tail-exponent}
\end{equation}
and there exists $A_\beta\in(0,\infty)$ such that, as $t\to\infty$,
\begin{equation}
 c(t)-q\sim A_\beta t^{-a},\qquad
 (-1)^nc^{(n)}(t)\sim A_\beta(a)_nt^{-a-n}\quad(n\ge1).
 \label{eq:main-stationary-tail}
\end{equation}
Here $(a)_n=a(a+1)\cdots(a+n-1)$, and $n$ is fixed.
\end{theorem}

The relation \eqref{eq:main-tail-exponent} is the exponent relation of mode-coupling theory for the stationary profile of the aging solution, with exponent parameter $x/2$ \cite[arXiv version, Equation~(IV.6)]{bckm1996}, \cite{kimlatz2001}. At the transition, $x=1$ and $a\approx0.395$ for every $p$; as $\beta\to\infty$, the right side of \eqref{eq:main-tail-exponent} tends to zero and $a$ tends to $1/2$. The critical power law, with an amplitude that depends on the dynamics at finite time differences, is part of the classical asymptotic analysis of mode-coupling theory \cite[Section~4, Equation~(6), pp.~880--881]{goetze1998}. In \cref{thm:stationary-critical-amplitude}, we prove it for the stationary profile selected by the CK solution, together with a positive finite amplitude and the asymptotics of every fixed derivative. In particular, the response $r=-c'$ decays like $aA_\beta t^{-a-1}$, and the approach to the plateau has no single exponential relaxation time. These exponents describe how the stationary profile approaches its plateau after the limit of large waiting time; they are not exponents of the aging regime. None of the geometric or Fisher-cost arguments uses this power law.

\section{Relation to previous work}\label{sec:intro-related}

We now compare our results and methods with previous work.

\paragraph{The physical predictions and their later refinements.}
The predictions tested by our results come from the asymptotic solution of Cugliandolo and Kurchan \cite{ck1993}, described in \cref{sec:intro-picture}. Their derivation assumes that the stationary and aging regimes separate, and it imposes an asymptotic form on the correlation and response in the aging regime. The time scale of the aging regime is not determined by this calculation, because the asymptotic equations admit reparametrizations of time \cite[arXiv version~1, Section~4]{ck1994}. In \cref{thm:main-aging}, we prove the stationary FDT, the marginal plateau, the energy, and the forgetting of fixed earlier times for the CK solution itself, without assuming that the regimes separate. \Cref{thm:main-response} identifies one weighted moment of the slow response with the ratio $x$, and \cref{thm:main-geometry} gives one side of the product relation between three times. The full slow fluctuation--dissipation relation and the time scale of aging remain open. The effective-temperature interpretation of $x$ was developed by Cugliandolo, Kurchan, and Peliti \cite{ckp1997}. Franz, M\'ezard, Parisi, and Peliti related the dynamical response to the equilibrium overlap distribution under additional hypotheses \cite{fmpp1998}, and they pointed out that the unperturbed spherical $p$-spin model does not satisfy this relation. We do not use equilibrium overlaps. Altieri, Biroli, and Cammarota developed a formulation of aging within dynamical mean-field theory in \cite{altieri2020}, and Chamon, Kennett, Castillo, and Cugliandolo developed the reparametrization invariance of the slow dynamics in \cite{chamon2002}. Ghimenti, Berthier, Kurchan, and van Wijland related accelerated glass dynamics to time-reparametrization softness, through numerical comparisons and an asymptotic analysis of a nonreciprocally coupled spherical $p$-spin model \cite[pp.~4--5 and SI Appendix, Section~7]{ghimenti2026}.

The picture has also been tested numerically and refined. Kim and Latz integrated the CK equations numerically over long times and compared the result with the predicted fluctuation--dissipation relation and with the crossover to the aging regime \cite{kimlatz2001}. Andreanov and Lef\`evre proposed a form of the time scale of aging by matching the stationary and aging regimes within a scaling expansion \cite{andreanovlefevre2006}. Folena, Franz, and Ricci-Tersenghi found that in zero-temperature quenches of mixed spherical models the long-time dynamics can depend on the initial condition and can reach energies below the threshold \cite[Sections~V--VII]{folena2020}. Folena and Zamponi tested weak ergodicity breaking numerically in mean-field spin glasses over long times, mainly in mixed models \cite{folenazamponi2023}, and Lang, Sachdev, and Diehl developed a numerical renormalization method that reaches longer integration times \cite{lang2025}. Citro and Ricci-Tersenghi obtained numerical evidence for strong ergodicity breaking, in which the correlation with the initial configuration stays positive at long times, in the spherical $3+4$ model quenched from near the dynamical transition to zero temperature \cite{citro2025}. The studies of mixed models and of other initial conditions indicate that the selection of the threshold in \cref{thm:main-aging} depends on our hypotheses.

\paragraph{Rigorous limits on fixed time horizons and the CK equations.}
The rigorous derivation of mean-field dynamical equations began with the large deviation methods on path space of Ben Arous and Guionnet for Langevin dynamics of spin glasses: a large deviation principle for asymmetric interactions \cite{bg1995} and a large deviation upper bound for symmetric interactions \cite{bg1997}, in regimes of short times or high temperature. Grunwald proved an analogous asymmetric large deviation principle and symmetric upper bound for Glauber dynamics, the spin-flip dynamics of Ising spin glasses, without restrictions on the fixed time horizon or temperature \cite{grunwald1996}. Guionnet obtained averaged propagation of chaos, that is, asymptotic independence of finitely many coordinates, and a quenched single-spin limit expressed as convergence in law over the disorder \cite{guionnet1997}. For spherical models, Ben Arous, Dembo, and Guionnet proved in \cite[Theorem~1.2]{bag2006} that the empirical correlation and integrated response of the soft-spherical Langevin dynamics converge on fixed time horizons to the solution of the CK equations. Dembo, Guionnet, and Mazza took the limit of these deterministic equations as the soft constraint becomes hard \cite[Proposition~1.1]{dgm2007}. In the same paper, they proved exponential relaxation and convergence to a stationary FDT solution at sufficiently high temperature \cite[Proposition~1.2 and Theorem~1.3]{dgm2007}, and they classified the solutions of the scalar memory equation by the largest admissible plateau and a stability inequality \cite[Proposition~1.4]{dgm2007}. They also showed that the CK solution has stationary FDT limits at every temperature, under the asymptotic conditions (6.5) and (6.6) of their paper \cite[Proposition~6.1]{dgm2007}. Sellke proved the convergence theorem for the hard-spherical diffusion itself \cite[Theorem~4.3]{sellke2024}, by coupling it with soft-spherical diffusions \cite[Lemmas~3.1 and~3.2]{sellke2024}. Guionnet and Mazza studied the long-time behavior of nonlinear Kraichnan equations, which are integral equations for response kernels with a prescribed covariance, in \cite{guionnetmazza2005}, using positive noncrossing expansions: series with nonnegative terms indexed by pairings of ordered points whose arcs do not cross. Dembo and Subag derived dynamical equations on fixed time horizons for initial conditions in prescribed bands near critical points \cite{ds2020} and for Gibbs distributed initial conditions \cite{ds2025}. The latter result includes the initialization at inverse temperature zero, which is the uniform start, together with the limit of a hard spherical constraint \cite[Theorem~1.8(a) and Proposition~1.10]{ds2025}. They also characterize stationary FDT dynamics and relaxation within equilibrium bands \cite[Proposition~1.14 and Corollary~1.18]{ds2025}. The equilibrium relaxation conclusions concern Gibbs starts at the ambient temperature, whereas our long-time selection follows a quench from the uniform law to inverse temperature $\beta>0$. Dembo, Lubetzky, and Zeitouni proved universality on fixed time horizons for Langevin-like dynamics with independent asymmetric interactions \cite[Theorem~1.1]{dlz2021}. For diffusions with linear random matrix interactions, Dembo and Gheissari proved universality of averaged observables for both symmetric and asymmetric couplings, using a stochastic Taylor expansion \cite[Theorem~1 and Corollary~3]{dembogheissari2021}. Surveys of this program are \cite{ba2002} and \cite[pp.~117--121]{guionnet2007}.

Rigorous dynamical mean-field theory has recently been extended to first-order methods with random data \cite{ccm2021} and to stochastic gradient descent \cite{gerbelot2024}, and its equations have been derived by a dynamical cavity method \cite{dandi2026}. Long-time results are known for Langevin dynamics in spiked matrix models, which are random matrices with a low-rank signal \cite{lss2023}. For Langevin diffusions with adaptive drift, Fan, Ko, Loureiro, Lu, and Shen proved convergence of the linear response on fixed horizons \cite{fan2025a}, and long-time results under an approximate time-translation invariance hypothesis, which they verify under a log-Sobolev inequality that is uniform in the dimension \cite[Theorems~2.5 and~2.9]{fan2025b}. For the orthogonally invariant Sherrington--Kirkpatrick model, Fan, Misiakiewicz, Wang, and Wen proved stationary FDT limits under rapid mixing on time scales that do not depend on the dimension and suitable confining curvature \cite[Lemma~4.1]{fan2026}. The models and hypotheses of these works differ from ours.

We take the results of \cite{bag2006,dgm2007,sellke2024} as input, through \cref{prop:ck-wellposed,prop:hard-bridge}, and prove long-time statements for the CK solution at every temperature. The main difference from \cite[Proposition~6.1]{dgm2007} is that our stationary limits (\cref{prop:stationary-clusters}) require no asymptotic hypothesis on the solution: entropy production supplies the needed control.

\paragraph{Aging for the two-spin model and on other time scales.}
Ben Arous, Dembo, and Guionnet proved aging for the soft-spherical model with $p=2$ in the order of limits used here \cite{bag2001}. Their analysis uses the spectral decomposition of the interaction matrix, which reduces the dynamics to coupled Ornstein--Uhlenbeck processes. This linear structure is absent for $p\ge3$, where a positive stationary plateau must satisfy the nonlinear marginality condition \eqref{eq:main-marginality}.

A different rigorous theory of aging concerns time scales that grow with the dimension, on which the dynamics is activated and is trapped in deep states for long times. Ben Arous, Bovier, and Gayrard proved aging for the Glauber dynamics of the random energy model below the critical temperature \cite{bbg2003a,bbg2003b}. Ben Arous and \v{C}ern\'y showed that the arcsine law of a stable subordinator is a universal aging scheme for trap models: under their hypotheses, the probability that the process occupies the same trap at two late times converges to a function of their ratio given by the generalized arcsine law \cite{bc2008}. For random-hopping-time dynamics, in which the jumps do not depend on the energy and the holding times do, Ben Arous, Bovier, and \v{C}ern\'y proved that Ising $p$-spin models with $p\ge3$ age like the random energy model on intermediate exponential time scales \cite{bbc2008}, and Ben Arous and G\"un obtained extremal aging, which is governed by the longest holding times, on subexponential time scales \cite{bg2012}. Bovier and Gayrard proved convergence of clock processes in random environments, and aging for random-hopping-time dynamics of the $p$-spin SK model \cite{boviergayrard2013}. For the Metropolis dynamics of the random energy model, \v{C}ern\'y and Wassmer proved convergence of the clock process and explained why aging of the two-time correlation does not follow directly from it \cite[Theorem~1.1 and Remark~1.2(3)]{cw2017}, and Gayrard proved two-time aging \cite{gayrard2019}. These results let the observation times grow with $N$ and describe the dynamics through clock processes, which record the total time spent in the states visited by the dynamics; this is the trap mechanism of \cite{bouchaud1992}. In the CK limit, the dimension tends to infinity first on fixed horizons, activation is invisible, and the aging of \cref{thm:main-aging} is a property of the deterministic CK solution. The differences between these orders of limits are discussed in the surveys \cite{ba2002,bc2006} and \cite[pp.~120--121]{guionnet2007}.

\paragraph{Landscape and algorithmic thresholds, shattering, and metastability.}
In physics, the threshold energy $\Eth=2\sqrt{(p-1)/p}$ is the energy of the marginal TAP states at zero temperature \cite[Equation~(9)]{cs1995}. Auffinger, Ben Arous, and \v{C}ern\'y used random matrix theory to compute the complexity of the critical points of the pure spherical $p$-spin Hamiltonian, that is, the exponential growth rate of their number. Up to the sign convention, $\Eth$ is the energy $E_\infty$ of their analysis \cite[Equation~(2.13)]{abc2013}. Auffinger and Ben Arous extended these computations to mixed models \cite{ab2013}. In pure models, Subag proved by a second moment computation that, for each fixed energy strictly between the threshold and the ground-state energy, the number of critical points with energy above it, divided by its expectation, converges to one in $L^2$, in our sign convention \cite[Theorem~1 and Corollary~2]{subag2017}. In optimization, $\Eth$ is the algorithmic threshold of pure models, that is, the largest energy that algorithms of a natural stable class can reach. Subag constructed an algorithm for spherical models, based on the Hessian, that reaches the energy $\int_0^1\sqrt{\xi''(u)}\,\dd u$, which equals $\Eth$ for $\xi(u)=u^p$ \cite{subag2021}. Montanari \cite{montanari2019} and El Alaoui, Montanari, and Sellke \cite{eams2021} constructed message passing algorithms for Ising models. These algorithms iteratively use the gradient of the Hamiltonian, and with probability tending to one, they reach energies arbitrarily close to the ground-state energy under a no-overlap-gap assumption on the Parisi order parameter: in \cite[Assumption~1]{montanari2019}, that the support of the Parisi measure is an interval $[0,q_*(\beta)]$ for all sufficiently large $\beta$, and in \cite[arXiv version~1, Assumption~2]{eams2021}, that the zero-temperature Parisi minimizer is strictly increasing on $[0,1)$. Huang and Sellke proved a matching hardness result. Consider algorithms that map the Gaussian disorder coefficients to points of the ball of radius $\sqrt N$ in $\R^N$, with Lipschitz constants, in the Euclidean metrics, bounded independently of $N$. Then the probability that the energy per coordinate of the output exceeds the algorithmic threshold by a fixed positive amount tends to zero as $N\to\infty$. They proved this for even models in \cite{hs2021}, and for multi-species spherical models, including odd pure models and external fields, in \cite[Theorem~1 and Corollary~1.8]{hs2023}. Sellke proved that at low temperature, Langevin dynamics reaches energies close to $\Eth$ on bounded time horizons \cite[Theorem~1.1]{sellke2024}. He also proved that, from an initial point independent of the disorder, it stays below $\Eth$ on each fixed time horizon with probability tending to one as $N\to\infty$ \cite[Theorem~1.7]{sellke2024}. Our energy formula agrees with these results when the zero-temperature limit is taken after the long-time limits (\cref{rem:low-temperature}).

For every stationary limit with positive plateau, we use $\Eth$ as an upper bound for the limiting energy of normalized prefix centers, which are the projections of late tree vectors described in \cref{sec:main-results} (\cref{prop:threshold-inputs}). For $p\in\{3,4\}$, we prove this bound with a deterministic inequality for unit vectors in a Hilbert space graded by ordered trees (\cref{thm:geometry-tree-ceiling,thm:geometry-quartic-tree-ceiling}). The inequality has a sharp constant, and the deficit from this constant controls the squared norm of the components of a vector at every fixed height, where the height of a tree is the largest distance from its root to a leaf. Branching constructions appear in the algorithms of Subag \cite[Remark~6]{subag2021} and in the overlap geometry of Huang and Sellke \cite[Remarks~1.2 and~1.3]{hs2021}; here we retain quantitative height defects inside the causal covariance realization. For general $p$, we use the theorem of Huang and Sellke \cite{hs2023} together with the coupling of Sellke \cite[Lemmas~3.1 and~3.2]{sellke2024}.

Sellke's work on the marginal stability of approximate ground states \cite{sellke2025} concerns a different part of the landscape, and we do not relate the stationary response to a Hessian. The finite-dimensional dynamics has also been studied directly. Gheissari and Jagannath estimated the spectral gap of spherical spin glass dynamics \cite{gj2019}, and Ben Arous, Gheissari, and Jagannath bounded the evolution of the energy and the gradient by differential inclusions \cite{bgj2020}. Ben Arous and Jagannath separated shattering, a decomposition into exponentially many well-separated and individually subdominant regions whose union has the full limiting free energy, from metastability, which concerns long residence times near a region before escape \cite{bj2024}. El Alaoui, Montanari, and Sellke proved shattering of the Gibbs measure of pure spherical models in a range of temperatures, and showed that Langevin dynamics started from the Gibbs measure keeps a large correlation with its initial point on exponentially long time scales \cite[Corollary~2.9]{shattering2025}. This plateau concerns equilibrium initial data and time scales that grow with $N$. The plateau in \cref{thm:main-aging} concerns a uniform initial point and the limit in which the dimension tends to infinity first.

\paragraph{Methods.}
We describe the origin of the main tools of the proofs; \cref{sec:proof-overview} explains how they fit together.

\emph{Entropy production and FDT.} Agarwal \cite{agarwal1972} treats response around nonthermal steady states under detailed balance; nonequilibrium diffusion response identities are given in \cite{bmw2009}. Cugliandolo, Dean, and Kurchan bounded violations of FDT by the entropy production in \cite[arXiv version~2, Equations~(12)--(14) and~(18)]{cdk1997}. We use the Bochner formula for a Laplacian with a gradient drift \cite[Proposition~3]{bakryemery1985} and the calculus of reversible diffusions \cite{bgl2014}. In \cref{sec:entropy-clusters}, we establish dimension-uniform bounds for the projected hard-spherical response and pass them to a single Fisher majorant for the CK solution, which controls the violation of FDT at all pairs of times $t>s>0$ and gives stationary FDT limits.

\emph{Effective processes and the Green identity.} The covariance identity of \cref{prop:green-covariance} has the variation-of-constants interpretation of the linear effective single-spin equation with white and colored Gaussian noise \cite[arXiv version~2, Section~9.5, Equation~(9.56), p.~55]{cugliandolo2004}. We prove the identity directly from the CK equations in \cref{sec:ck-structure}.

\emph{Expansions with nonnegative terms.} Positive noncrossing expansions of the response equation appear in the work of Guionnet and Mazza \cite[Section~2.1]{guionnetmazza2005} and of Dembo, Guionnet, and Mazza \cite[Section~2]{dgm2007}. In \cref{sec:geometry-core}, we expand the covariance instead. We organize the expansion by ordered trees whose leaves record either the initial condition or a time of the thermal noise. This realizes $C(t,s)$ as the inner product of unit vectors in one deterministic Hilbert space, and the projection onto trees with early leaves plays the role of conditioning on the early noise.

\emph{Independent futures.} Trajectories with a common past and independent future noise were used by Cugliandolo and Dean \cite[arXiv version~1, Section~II]{cugliandolodean1995} and by Barrat, Burioni, and M\'ezard, who classified aging by the overlaps of such trajectories \cite[Section~2]{bbm1996}. The overlap of two independent continuations of lengths $a$ and $b$ equals the correlation at the reflected lag $a+b$ in equilibrium \cite[Section~2, Equation~(9)]{bbm1996}. Franz and Virasoro proposed averaging cloned futures to estimate the magnetization of the region of configurations that the dynamics explores on intermediate time scales \cite[arXiv version~3, Section~III]{franzvirasoro2000}, and noise-averaged magnetizations and their relation to the threshold are central to the dynamical TAP approach of Biroli \cite[arXiv version~2, Section~III.B]{biroli1999}. Time-averaged magnetizations whose normalized squared norm tends to $q$ also appear in \cite[arXiv version~1, Section~5.1, Equation~(5.1.9)]{ck1994}. Our prefix centers are deterministic counterparts of these averages. In \cref{sel:prefix-reflection}, we prove the reflected formula for the stationary limits of the CK solution: the older memory cancels exactly in a closed integral equation for the times after the earlier time, before any limit is taken. For general $p$, \cref{app:general-selection} uses finitely many branched future trajectories of the diffusion.

\emph{Complete monotonicity.} The reflected kernel $(a,b)\mapsto c(a+b)$ is positive semidefinite, and together with the boundedness and continuity of $c$, this implies that $c$ is completely monotone, by the Bernstein--Widder theorem and its semigroup form \cite{widder1941,bcr1984}. Franosch and Voigtmann proved complete monotonicity for Brownian mode-coupling equations of mixtures by a different route, an iteration argument \cite[arXiv version~1, Section~III]{fv2002}.

\emph{Gaussian conditioning.} In \cref{app:general-selection}, we prove forgetting for general $p$ by revealing the disorder in stages, through the gradients of the Hamiltonian at finitely many points chosen from the history, and by applying the external-field threshold of \cite[Theorem~1 and Corollary~1.8]{hs2023} at each stage. An Euler approximation of an earlier state lies in the span of the initial point, finitely many Brownian increments, and the computed gradients, and we condition on the Gaussian coefficients needed to determine the gradients on this span. Gaussian conditioning is classical \cite{janson1997}. Adaptive Gaussian conditioning is used by Bayati and Montanari for approximate message passing \cite[arXiv version~4, Section~3.8]{bayatimontanari2011} and by Bolthausen for the iterative construction of solutions of the TAP equations \cite[Section~3]{bolthausen2014}. El Alaoui, Montanari, and Sellke conditioned the disorder tensor on adaptive gradient observations in their analysis of message passing for spin glasses \cite[arXiv version~1, Appendix~A]{eams2021}, and Huang and Sellke extended this conditioning to multi-species spherical models \cite[Appendix~A]{hs2024amp}. We adapt the method to a fixed finite history of the Langevin dynamics.

\emph{Mode-coupling and renewal tools.} The scalar memory equation of a stationary profile belongs to the class studied in mode-coupling theory \cite[Sections~3.1--3.2]{goetzesjogren1992} and \cite[Section~4.2.1, Equations~(4.15)--(4.17), pp.~192--193]{goetze2009}. The structural theory of these equations is developed in \cite{gs1995} and \cite{fv2002}, including the characterization of the plateau \cite[arXiv version~1, Section~IV]{fv2002}, and Sj\"ogren interpreted them through regenerative processes, which restart at random times \cite{sjogren2003}. Our proof of \cref{thm:main-stationary-tail} in \cref{sec:stationary-relaxation} combines comparison principles, a weighted Schur inequality for the convolution in the scalar equation, a factorization of gamma variables, which turns a quadratic balance into a renewal equation, the key renewal theorem \cite{feller1971}, and the monotone density theorem \cite[Section~1.7.3]{bgt1987}.

\emph{The three-time inequality.} The product relation for three times is part of the physical analysis of Cugliandolo and Kurchan \cite[arXiv version~1, Section~5.1, Equation~(5.1.8)]{ck1994}. We prove a uniform one-sided inequality for the CK solution in \cref{thm:geometry-actual-product}, from the constrained geometry of tree vectors whose energies approach the threshold (\cref{thm:geometry-threshold-product}). We use it for the bound on the Fisher cost in \cref{thm:main-fisher-cost}.

\paragraph{Concurrent work.}
Several recent works study related questions with different models, initial conditions, or hypotheses. Dembo and Subag analyzed dynamics started from Gibbs measures, including limits on fixed time horizons and the stationary FDT dynamics and equilibrium band relaxation described above \cite{ds2025}. Their theorem includes the initialization at inverse temperature $\beta_0=0$, hence the uniform start \cite[Theorem~1.8(a) and Proposition~1.10]{ds2025}; for initialization from a low-temperature Gibbs measure, \cite[Theorem~1.8(b)]{ds2025} assumes additional conditions on the mixture and on the equilibrium Parisi measure (genericity and strict one-step replica symmetry breaking, in the sense of \cite[Section~1]{ds2025}). The works \cite{fan2025a,fan2025b,fan2026} on dynamical mean-field equations, described above, concern other models; their long-time results use approximate time-translation invariance or mixing hypotheses. Lang, Citro, Leuzzi, and Ricci-Tersenghi computed the aging phase diagram and the asymptotic energies of mixed spherical models at zero temperature, within a formal asymptotic analysis \cite{lang2026}. Cipolloni and Cugliandolo studied coupled spherical $p$-spin systems, including their aging regimes and fluctuation--dissipation relations \cite{cipolloni2026}. Our contribution relative to these works is the set of long-time conclusions for the uniformly initialized CK solution of pure models: the selection of the plateau and the energy, the forgetting of fixed earlier times, and the results on the response.

\section{Scope and open problems}\label{sec:intro-scope}

Our conclusions concern the CK solution started from a uniform point, not an arbitrary solution of a reduced asymptotic system, and they apply to Langevin dynamics only through the limit on fixed time horizons of \cref{prop:hard-bridge}. For every $p\ge3$, the stationary limits, the convergence of the energy, the forgetting of fixed earlier times, and the convergence of the response over all earlier times hold at every positive temperature. The identification of the plateau and the energy, correlation aging, and the stationary power law hold for $\beta\ge\bd$, and the weighted ratio and the persistent violation of FDT hold for $\beta>\bd$. For the cubic model, the three-time inequality and the susceptibility bound hold for $\beta\ge\bd$, while the positive Fisher cost requires $\beta>\bd$.

Several parts of the physical picture remain open, and we formulate them in \cref{sec:discussion}:
\begin{itemize}
\item the full slow fluctuation--dissipation relation, of which \eqref{eq:main-weighted-ratio} is one weighted moment (\cref{sec:discussion-fdr});
\item an upper bound matching \eqref{eq:main-susceptibility} for the susceptibility (\cref{sec:discussion-susceptibility});
\item the time scale of aging, which the three-time inequality and the Fisher cost constrain but do not determine (\cref{sec:discussion-clock});
\item the selection of the plateau for $\beta_{\mathrm m}\le\beta<\bd$ (\cref{sec:discussion-below});
\item the behavior on time scales that grow with the dimension, and at zero temperature (\cref{sec:discussion-limits}).
\end{itemize}

\section{Organization of the paper}\label{sec:intro-organization}

\Cref{sec:proof-overview} gives an overview of the proofs. They have three main ingredients. From entropy production, we obtain stationary FDT limits on bounded late windows. The representation of the correlation by vectors in a Hilbert space indexed by trees separates the contributions of the earlier and the later noise. An energy bound for the earlier contributions selects the marginal plateau and shows that fixed earlier configurations are forgotten. The rest of the paper has five parts.

\Cref{part:foundations} contains the foundations. \Cref{sec:model} states the model, the conventions, the CK equations, and the published inputs, including the limit theorem on fixed time horizons. In \cref{sec:ck-structure}, we prove the energy identity, the Green covariance identity, and estimates that are uniform in time. In \cref{sec:entropy-clusters}, we use entropy production to construct the Fisher majorant and to show that all limits of the CK solution on late time windows are stationary and satisfy FDT.

\Cref{part:selection} proves \cref{thm:main-aging}. \Cref{sec:scalar} derives the scalar equation of a stationary limit and its constraints. \Cref{sec:geometry-core} realizes $C$ as the Gram kernel of unit vectors in a Hilbert space graded by trees, and proves the reflection identity and complete monotonicity. \Cref{sec:stationary-selection} states the two threshold inputs (\cref{prop:threshold-inputs}) and deduces from them the selection of the plateau, the forgetting of earlier times, and correlation aging. \Cref{sec:tree-ceilings} proves the threshold inputs for $p\in\{3,4\}$ with deterministic inequalities in the tree space, and \cref{app:general-selection} proves them for every $p\ge3$.

\Cref{part:response} proves \cref{thm:main-response}. In \cref{resp:chapter}, we prove the convergence of the response and of the derivatives of the correlation over all earlier times, and the convergence of the thermal term. In \cref{resp:integrated-chapter}, we prove the weighted ratio and the persistent violation of FDT.

\Cref{part:geometry} concerns the cubic model and proves \cref{thm:main-geometry,thm:main-fisher-cost}. \Cref{sec:geometry-results} studies tree vectors near the threshold, \cref{sec:three-time} transfers their geometry to the CK correlations and proves the susceptibility bound, and \cref{resp:fisher-chapter} proves the bound on the Fisher cost.

\Cref{part:relaxation} contains \cref{sec:stationary-relaxation}, which proves \cref{thm:main-stationary-tail}, and \cref{sec:discussion}, which discusses the open problems. \Cref{sec:stationary-relaxation} can be read once \cref{part:selection} is available.

\Cref{tab:main-results-map} lists, for each main theorem, where it is proved and the main results that the proof uses.

\begin{table}[ht]
\centering
\small
\renewcommand{\arraystretch}{1.3}
\begin{tabular}{@{}>{\raggedright\arraybackslash}p{0.15\textwidth}>{\raggedright\arraybackslash}p{0.27\textwidth}>{\raggedright\arraybackslash}p{0.5\textwidth}@{}}
\toprule
Result & Proved in & Main inputs\\
\midrule
\Cref{thm:main-aging} & \Cref{sec:stationary-selection}; the proof is assembled at the end of \cref{sec:aging} & \Cref{prop:hard-bridge}; \cref{prop:stationary-clusters}; \cref{prop:scalar-clusters}; \cref{sel:prefix-reflection,sel:cluster-complete-monotonicity}; \cref{prop:threshold-inputs}, proved in \cref{sec:tree-ceilings} for $p=3,4$ and in \cref{app:general-selection} for all $p$\\
\Cref{thm:main-response} & \Cref{thm:fisher-clock}; \cref{resp:strong-row,resp:uniform-row,resp:backward-fdt}; \cref{resp:weighted-ratio,resp:persistent-defect} & \Cref{thm:main-aging}; \cref{prop:green-covariance}; \cref{sel:prefix-reflection}; \cref{thm:geometry-tree-ceiling} ($p=3$) or \cref{resp:appendix-p} (all $p$); \cref{resp:transport}\\
\eqref{eq:main-thermal} & \Cref{resp:thermal} & \Cref{sel:cluster-complete-monotonicity}; \cref{resp:strong-row}\\
\Cref{thm:main-geometry} & \Cref{thm:geometry-actual-product,thm:geometry-susceptibility-lower} & \Cref{thm:geometry-tree-ceiling}; \cref{sel:threshold-centers}; \cref{thm:geometry-universal-root,thm:geometry-threshold-product}; \cref{resp:strong-row,resp:uniform-row,resp:forward-derivative}\\
\Cref{thm:main-fisher-cost} & \Cref{resp:fisher-cost} & \Cref{prop:backward-response}; \cref{thm:fisher-clock}; \cref{resp:uniform-row,resp:forward-derivative}; \cref{thm:geometry-actual-product}\\
\Cref{thm:main-stationary-tail} & \Cref{thm:stationary-critical-amplitude} & \Cref{prop:scalar-clusters}; \cref{sel:cluster-complete-monotonicity}; \cref{sel:selected-sector}\\
\bottomrule
\end{tabular}
\caption{The main results, the sections and results in which they are proved, and the main inputs of their proofs.}
\label{tab:main-results-map}
\end{table}
